
\begin{table}[H]
    \centering
    \caption{Interrupt Vector Entry Address}
    \label{table:Interrupt}
    \begin{tabular}{|p{2.5cm}|p{6.5cm}|p{6cm}|}
        \hline
        \rowcolor{lightgray}
        Priority level &
        \makecell[l]{PIDCTRL\_INTERRUPT\_ENABLE\_REG bit\\controlling interrupt identification} &
        Interrupt vector entry address \\
        \hline
        Level 1 & 1 & \hyperref[regdesc:PIDCTRLINTERRUPTADDR1REG]{PIDCTRL\_INTERRUPT\_ADDR\_1\_REG}\\
        \hline
        Level 2 & 2 & \hyperref[regdesc:PIDCTRLINTERRUPTADDR2REG]{PIDCTRL\_INTERRUPT\_ADDR\_2\_REG} \\
        \hline
        Level 3 & 3 & \hyperref[regdesc:PIDCTRLINTERRUPTADDR3REG]{PIDCTRL\_INTERRUPT\_ADDR\_3\_REG}\\
        \hline
        Level 4 & 4 & \hyperref[regdesc:PIDCTRLINTERRUPTADDR4REG]{PIDCTRL\_INTERRUPT\_ADDR\_4\_REG}\\
        \hline
        Level 5 & 5 & \hyperref[regdesc:PIDCTRLINTERRUPTADDR5REG]{PIDCTRL\_INTERRUPT\_ADDR\_5\_REG}\\
        \hline
        Level 6 ( Debug ) & 6 & \hyperref[regdesc:PIDCTRLINTERRUPTADDR6REG]{PIDCTRL\_INTERRUPT\_ADDR\_6\_REG}\\
        \hline
        NMI & 7 & \hyperref[regdesc:PIDCTRLINTERRUPTADDR7REG]{PIDCTRL\_INTERRUPT\_ADDR\_7\_REG}\\
        \hline
    \end{tabular}
\end{table}
